\introduction{一切结构的尺度}
《魔法少女小圆》（以下简称《魔圆》，以区别于主角鹿目圆\footnote{%
    日语习惯将姓氏置于名字之前（如鹿目圆），但《魔圆》官方英译采用名在前、姓在后的顺序。
    为避免混淆，本书沿用剧中顺序。
    当涉及其他角色或真实人物时，为求统一，亦采用此顺序。
}，简称“小圆\footnote{%
    译注：
    原书按英语习惯名前姓后；
    译文统一按中文通行方式采用姓前名后。
}”）内涵极为丰富。
显而易见，它是由动画公司 SHAFT 打造的一部魔法少女类型动画。
但除此之外，它不仅是对该类型的批判与致敬，也是一位迷失于黑暗中的作者寻求自我救赎的尝试，还可视为从佛教视角对《浮士德》的重述。
它更是一部关于对立与统一的深刻之作。
它既从女性主义视角对同意、特权及性别角色展开严肃拷问，同时又（不无矛盾地）对未成年女性的身体施以带有性别歧视色彩的物化。
它所书写的是牺牲与自私、绝望与希望、爱与死亡、衰败与新生，并最终揭示出：
所有这些看似水火不容的对立面，实则浑然一体，并无二致。

然而，归根结底，它讲述的是人的故事\footnote{%
    译注：
    导言标题“一切结构的尺度（The Measure of All Structures）”可能化用自普罗泰戈拉名言：“人是万物的尺度（Man is the measure of all things）。”
}。
故事中的人物作出选择（往往是错误的选择），承受苦难、奋力抗争，也经历成功与失败。
角色之间或彼此相爱，或彼此憎恨，或漠不关心，有时甚至三者并存。

本书将探讨诸多宏大主题、哲学命题与象征符号。
读者请谨记：
一切故事皆关乎人\footnote{%
    “摹仿者所摹仿的对象是行动中的人。”
    亚里士多德的《诗学》第一部分，将戏剧与叙事纳入诗学范畴，其所谓“摹仿的对象”即今之“作品主题”。
}\cite{aristotle2009poetics}。
\paragraph{为什么研究《魔圆》？}

仅用十二集每集二十二分钟的电视动画，再加一部剧场版（虽有三部剧场版，前两部实为总集篇，改动甚微），《魔圆》便容纳了惊人的故事体量：
角色塑造饱满，主题深邃复杂，情感冲击强烈——深沉的悲伤、惊惧、愤怒与喜悦交织。
分析这部作品总是令人获益匪浅，每每我深入考察之时，总能发掘出新意蕴。
尤其是这部作品本身便具有鲜明的后现代色彩，因此用后现代视角来剖析它，相得益彰。
\paragraph{后——什么来着？}

后现代主义很难界定，但我还是想试着说明，至少说明本书所说的后现代主义究竟是什么。
从哲学上说，它是在现代主义之后又迈出了一步。

现代主义的核心洞见在于：
符号本质上具有任意性。
换言之，能指（表意之言）与所指\footnote{%
    译注：
    语言学术语，能指表示符号的形式，如语音、字形。
    所指表示符号的内容、概念。
}（被指之物）间并无天然纽带，其关联仅存于观者意识之中。
例如，红色六边形意指“停”，并非红、六边、停止三者存在逻辑关联，实则因有人规定此形此色表示“停”，并获得了众人认可。
此即社会建构主义内核：
我们理解世界的概念与符号，皆由社会环境与关系构筑而成。
故符号意义如社会关系般流动不居，并无固定客观标准\cite{burr1995social-constructionism}。

在后现代主义视域下，一切艺术（及语言）皆为符号系统。
因此，任何作品的意义在某种程度上皆具任意性，乃社会建构之物，非逻辑必然之果。
传统艺术倚赖约定俗成的文化符号，如语言、具象图像（比方说用苹果的图片代指苹果），来尽量避免任意解读。
此类符号如同“引导轨道”，帮助受众以熟悉的符号为起点接触艺术，并由此逐步构建更多意义。
现代主义大体以拒斥传统以及幻灭感、解体感为特征，其艺术常摒弃部分乃至全部“引导轨道”，激进地质疑“艺术具有——或应当具有——非任意的意义”这一观念。
例如，詹姆斯·乔伊斯\footnote{%
    译注：
    詹姆斯·乔伊斯（1882 年 2 月 2 日—1941 年 1 月 13 日），爱尔兰作家和诗人，现代主义文学的奠基者之一。
    代表作有《都柏林人》《尤利西斯》等。
}的作品以意识流打破句法常规；
蒙德里安\footnote{%
    译注：
    彼得·蒙德里安（1872 年 3 月 7 日—1944 年 2 月 1 日），荷兰画家、风格派运动领袖。
    代表作有《灰色的树》《纽约市一号》《红、黄、蓝的构成》等。
}的画作以抽象色块取代具象物体。
然而矛盾的是，现代主义常试图重建其感知中已失落的秩序\cite{murfin1997modernism}。
例如，乔伊斯将平凡的一天写成史诗般的抗争，蒙德里安则将抽象色块纳入规整网格中\footnote{%
    译注：
    这两个例子分别指代乔伊斯的《尤利西斯》和蒙德里安的《红、黄、蓝的构成》。
}。

后现代主义反其道而行，反对“社会建构即无意义”之论。
现代主义断言：
“能指与所指的关系具有任意性，故万物皆无意义。”
后现代主义主张：
“能指与所指的关系具有任意性，故万物皆可被自由地赋予意义。”
因此，后现代作品无需复现虚妄的“失落秩序”。
出于同样的原因，后现代主义更拥抱流行文化作为艺术表达媒介\footnote{%
    不过，Murfin 与 Ray 的定义着重强调后现代主义的消极面与实验性，淡化了伊哈布·哈桑视为其核心特征的多元主义。
}\cite{murfin1997postmodernism}。
现代主义摒弃传统以求新秩序；
后现代主义则糅合传统与新异，追寻奇趣之境\footnote{%
    伊哈布·哈桑列出了十一项彼此交叠的判准。本书给出的定义主要取自其中的“建构性”“去经典化”“反讽”“狂欢性”与“杂交”。
}\cite{hassan1987pluralism}。

如果说现代艺术是拆除“引导轨道”，那么后现代艺术则是刻意凸显“引导轨道”。
其核心手法是去语境化与再语境化：
将熟悉的能指剥离原境，植入异质语境，如混搭多类型元素、仿拟人们熟悉的作品、打破第四面墙\cite{hassan1987pluralism}。
由此引发的错位感，常被用于制造幽默或恐怖效果，或激发思辨（三者亦可并存）。

后现代主义有“积极”一脉，崇尚游戏与自由；
亦有“消极”一支，强调混乱、犬儒与绝望\cite{rosenau1996affirmatives-skeptics}。
《魔圆》大体属于后者，并大量运用了后现代手法叙事。
譬如（下章详述）“魔女”以异质侵凌的艺术风格呈现，强行取代动画一贯的画风。
\paragraph{本书结构}

本书分三部分。
第一部分为“电视动画”，十二章对应十二集，每章标题截取捷克剧作家、活动家、政治家瓦茨拉夫·哈维尔\footnote{%
    译注：
    瓦茨拉夫·哈维尔（1936 年 10 月 5 日—2011 年 12 月 18 日），首任捷克共和国总统。
}名言片段，括号内注明了该集动画标题。
选用哈维尔语录的用意，随阅读推进自明。
书名亦出自其言：
``或许，绝望恰是孕育人类希望的土壤；
未经荒诞，难觅生命真义。 % 此句全文多次出现，若需修改，请同步。
"\footnote{%
    译注：
    原文为``Perhaps hopelessness is the very soil that nourishes human hope;
    perhaps one could never find sense in life without first experiencing its absurdity.''
}\cite{shipley2003havel}。

第一部分之后设有间奏一“至福一瞬”，探讨本剧对《浮士德》叙事元素的运用。

第二部分为“外传漫画”，篇幅最短，三章分述电视动画播映期间及之后推出的三部漫画衍生作。
各章标题均引自电视动画中的台词。

第二部分之后是间奏二“牛奶之尸”，解析剧场版《魔法少女小圆【新篇】叛逆的物语》（以下简称《叛逆》）中奶酪的意义，以及炼金术中的“腐化”概念。

第三部分为“叛逆的物语”，共七章，聚焦剧场版《叛逆》。
本书只讨论这一部，是因为前两部剧场版只是电视动画的总集篇，新画面寥寥，更谈不上新增剧情，而《叛逆》实为电视动画的正统续篇\footnote{%
    影片初段有些模棱两可，看似是一条“平行时间线”，类似第二部分外传漫画的设定。
    然而上映的宣发明示：
    前两部为“回顾”，此部乃“延续”。
    角色保有最终话的记忆（在焰最后一次重置时间线与瓦尔普吉斯之夜决战之后）。
    例如剧中麻美与焰谈及“魔兽”；
    丘比也记得焰告诉过它小圆为创造新系统而牺牲的事；
    以及小圆最终以“圆环之理”的形态出现在结尾。
}。
七章分别讨论片名中的“叛逆”可能针对的七种对象，故标题皆作“反抗\underline{\makebox[1cm]{}}”。

本书是一部分析性著作，而非观剧指南。书中会在与分析相关时介绍情节、角色及制作细节，但总体上预设读者已经熟悉原作。
成书时，电视动画在美国可通过 DVD、Blu-ray 以及 Netflix、Crunchyroll 等流媒体平台观看。
相关漫画的英文译本也已在美国面世，或正在陆续出版。
《叛逆》仅曾在影院短暂上映，较难获取，但可以购买带英文字幕的日版 Blu-ray，配有英文字幕和英语配音的美版 Blu-ray 则定于 2015 年 4 月全面发行\footnote{%
    译注：
    国内可以通过哔哩哔哩视频网站观看。
    若需下载，推荐 VCB-Studio 压制组。
}。
基于这一预设，本书不再逐一标注剧中情节的出处，只在引用《魔圆》系列以外的文献时提供引注。

最后特别说明：
本书中有几章（尤其是第二和第九章）从女性主义角度分析了本剧的若干方面及魔法少女类型。
相关分析所援引理论植根西方，依托欧洲文化历史与预设，故其适用于西方观众，却未必契合日本文化——即该动画的创作文化和主要受众。
此部分论述，当视为西方视角下的解读，其结论未必能推广至日本观众的观看经验。
\paragraph{致谢}

本书成稿仰赖多方襄助。
首谢挚友兼封面设计师 Viga Gadson，引我初识《魔圆》。
次谢 Unnoun，他不断将我博客上的电视动画剧评转发至各个论坛，使这些文章成为我作品中访问量最高的一批，也促使我写下更多《魔圆》评论，最终汇集成本书。
三谢责编 Kit Paige，若这本书还有一点像样学术著作的门面，那都多亏了她；
若有疏漏，皆因我固执，非其失察。
四谢 MarkWatches.net 上的 Mark Oshiro 与 SFDebris.com 上的 Chuck Sonnenburg，二人风格迥异的评论使我得以跳出自身观点，获得必要的批评距离。
最后铭谢博客 JedABlue.com\footnote{%
    译注：
    原博客因故不可访，今存于 Internet Archive。
}上的读者。
诸位对大部分章节初稿提出的真知灼见，帮助我发现了其中的缺陷、需要澄清之处，以及我未曾考虑的反面观点。
尤谢 Alan Jacobs Richardson（01d55）：
我与他围绕《叛逆》和 E. T. A. 霍夫曼\footnote{%
    译注：
    E. T. A. 霍夫曼（1776 年 1 月 24 日—1822 年 6 月 25 日），德国作家、作曲家。
    其作品多具神秘、怪诞色彩，常以夸张手法讽刺、揭露现实。
    作品对人际关系异化的描写，以及自由联想、内心独白、夸张荒诞的表现和多层次结构等手法，都与后来的现代主义文学有着深厚渊源。
}《胡桃夹子与老鼠王》之间的关联展开了长篇讨论\footnote{%
    讨论见博客。
}\cite{blue2014nutcracker}，催生了第二十和二十一章的大量内容。
